\documentclass[11pt,letterpaper]{article}
\PassOptionsToPackage{hyphens}{url}
\usepackage[margin=.83in,top=.8in,bottom=.8in,headheight=15pt]{geometry}
\usepackage[T1]{fontenc}
\usepackage{lmodern,amsmath,amssymb,mathtools,bm,booktabs,array,graphicx,tikz}
\usetikzlibrary{arrows.meta,positioning,calc,decorations.pathreplacing,patterns}
\definecolor{ink}{HTML}{17344D}\definecolor{accent}{HTML}{087E8B}\definecolor{pale}{HTML}{EDF5F5}\definecolor{gold}{HTML}{AD6B19}
\providecommand{\dd}{\,\mathrm{d}}\providecommand{\R}{\mathbb{R}}
\providecommand{\E}{\mathbb{E}}\providecommand{\Var}{\operatorname{Var}}
\providecommand{\Cov}{\operatorname{Cov}}\providecommand{\diag}{\operatorname{diag}}
\providecommand{\norm}[1]{\left\lVert#1\right\rVert}
\providecommand{\abs}[1]{\left|#1\right|}
\usepackage{microtype,hyperref,fancyhdr,enumitem,amsthm,sectsty,needspace}
\allsectionsfont{\raggedright}
\hypersetup{colorlinks=true,linkcolor=ink,urlcolor=accent,pdfauthor={Zan Ahmad},pdftitle={Membrane voltage and a traveling activation / Review}}
\setlength{\parindent}{0pt}\setlength{\parskip}{6pt}\linespread{1.07}
\setlist{itemsep=3pt,topsep=4pt,leftmargin=*}
\emergencystretch=2em\clubpenalty=10000\widowpenalty=10000
\pagestyle{fancy}\fancyhf{}\fancyhead[L]{\small Mathematical Cardiology}\fancyhead[R]{\small Lecture 4}
\fancyfoot[L]{\footnotesize Zan Ahmad}\fancyfoot[R]{\footnotesize\thepage}\renewcommand{\headrulewidth}{.3pt}
\newtheorem{proposition}{Proposition}\newtheorem{theorem}{Theorem}\newtheorem{lemma}{Lemma}
\theoremstyle{definition}\newtheorem{definition}{Definition}\newtheorem{example}{Example}
\numberwithin{equation}{section}
\begin{document}
\section*{Membrane voltage and a traveling activation --- review sheet}
Use this after attempting the worked examples and practice. Explain each claim in words before calculating.

\section*{Key ideas}
\begin{itemize}
\item Membrane voltage, activation maps and ECG recordings are different objects.
\item Cell-specific currents and recovery matter; a neuronal analogy is not a cardiac model.
\item Constant-speed arrival time is distance divided by speed.
\item A prescribed delay translates a shape but does not solve tissue dynamics.
\item A 3D rendering does not establish 3D physiological detail.
\end{itemize}
\section*{Common errors and repairs}
\textbf{The electrical wave is blood moving through the tissue.} It is sequential excitation of coupled cells, not bulk blood transport.

\textbf{A delayed local trace is automatically an ECG lead.} An ECG requires a different observation model involving distributed activity and the body.

\section*{Retrieve without your notes}
\begin{itemize}
\item Convert 0.4 mm/ms to m/s.
\item What evidence could assess the constant-speed assumption?
\end{itemize}
\section*{Readiness check}
\begin{itemize}
\item Distinguish membrane voltage at one cell from a spatial activation pattern and an ECG.
\item Explain qualitatively how regulated ion movement supports excitation and recovery.
\item Calculate arrival time from distance and a specified conduction speed with consistent units.
\item Interpret a space-time plot and a 3D activation view without claiming that a prescribed wave solves tissue dynamics.
\end{itemize}
\end{document}
